\chapter{Evaluation Overview}
\label{ch:eval-overview}

% TODO: Once this chapter is integrated, remove the hardware descriptions from
%       evaluation_static.tex (\autoref{sec:eval:static:setup}),
%       evaluation_dynamic.tex (\autoref{sec:eval:dynamic:setup}), and
%       evaluation_property.tex (\autoref{sec:eval:property:setup}).
%       Also remove the datasets table from evaluation_static.tex (it now
%       lives here as \autoref{tab:datasets}).

\project{} claims to serve the full spectrum of graph workloads---from
whole-graph analytics to fine-grained transactional queries---on a single
persistent representation, without data movement or ETL\@.
Most existing systems optimize for one operating point: in-memory graph
processing systems build purpose-built layouts for analytics but require
offline preprocessing and cannot handle mutations; dynamic graph containers
sustain high ingest throughput but sacrifice analytical performance or
persistence; graph databases handle OLTP queries well but collapse on iterative
analytics.
If \project{} merely matched one class of system while failing on another, it
would reproduce the specialization problem it is designed to solve.
The evaluation must therefore test \project{} on each class of workload,
measuring it against the systems that specialize in that class.

We organize the evaluation into four workload classes, each targeting a
distinct operating regime:

\begin{enumerate}
  \item \textbf{In-memory graph analytics} (\autoref{sec:eval:inmem}).
        Can \project{} match specialized graph processing systems that build
        purpose-built in-memory layouts?
        We compare against GAPBS, Galois, Ligra, and Gemini on PageRank,
        connected components, BFS, and betweenness centrality, measuring both
        kernel execution time and the preprocessing overhead that each system
        requires before it can run algorithms.

  \item \textbf{Out-of-core graph analytics} (\autoref{sec:eval:ooc}).
        Does \project{} remain competitive when the graph exceeds available
        memory?
        We compare against GraphChi, X-Stream, and Blaze under constrained
        memory budgets, measuring kernel execution time, preprocessing
        overhead, and I/O bandwidth utilization.

  \item \textbf{Dynamic structural graphs} (\autoref{sec:eval:dynamic}).
        Can \project{} sustain mutation throughput and deliver analytics on a
        live, changing graph?
        We compare against Teseo, LiveGraph, and GraphOne on insertion
        scalability, mixed insert-delete workloads, and post-ingestion
        analytics (PageRank, BFS, weakly connected components).

  \item \textbf{Property graph queries} (\autoref{sec:eval:property}).
        Can \project{}'s property storage layer compete with mature graph
        databases on transactional queries?
        We compare against Neo4j, ArangoDB, NeuG, and others on individual
        structural and property CRUD operations, composite LDBC SNB queries,
        and full-graph algorithms.
\end{enumerate}

\noindent
The first two classes form \autoref{sec:eval:static} (static graph analytics);
the third forms \autoref{sec:eval:dynamic:chapter} (dynamic structural graphs)
and the fourth forms \autoref{sec:eval:property:chapter} (property graph
workloads).
The sections below describe the experimental environment, datasets, and
comparison systems shared across all four evaluations.

%% --------------------------------------------------------------------------
\section{Experimental Environment}
\label{sec:eval:env}

\autoref{tab:eval-hardware} summarizes the two machines used across all
experiments.
Machine~A runs the structural graph experiments (static analytics and dynamic
graph operations), where parallelism across many cores and large memory
capacity test scalability.
Machine~B runs the property graph experiments (microbenchmarks, LDBC SNB
queries, and graph algorithm comparisons against graph databases), where
single-threaded query latency matters more than core count.
Both machines run Ubuntu~22.04~LTS with Linux~6.8.0; we compile all
\project{} code and comparison systems with GCC~13.3.0 using
\texttt{-O3~-march=native~-mtune=native}.

%%%%%%%%%%%%%%%%%%%%%%%%%%%%%%%%%%%%%%%%%%%%%%%%%%%%%%%%%%%%%%%%%%%%%%%%%%%%%%%%
\begin{table}[h]
  \centering
  % \footnotesize
  \begin{tabular*}{\columnwidth}{@{\extracolsep{\fill}}lll@{}}
    \toprule
    & Machine~A & Machine~B \\
    \midrule
    Processor & AMD EPYC 7643 ($\times$2) & Intel Xeon W-2275 \\
    Clock     & 2.30\,GHz                 & 3.30\,GHz \\
    Cores / threads & 96 / 192            & 14 / 28 \\
    Memory    & 2\,TB                     & 128\,GB \\
    Storage   & 2\,TB NAND Flash SSD      & 1\,TB NVMe SSD \\
    \bottomrule
  \end{tabular*}
  \caption{Hardware configurations.  Hyperthreading enabled on both machines.}
  \label{tab:eval-hardware}
\end{table}
%%%%%%%%%%%%%%%%%%%%%%%%%%%%%%%%%%%%%%%%%%%%%%%%%%%%%%%%%%%%%%%%%%%%%%%%%%%%%%%%

Each comparison system runs inside a Docker container.
For the static and dynamic experiments on Machine~A, we restrict each container
to one NUMA node (96~threads) with local memory; the out-of-core experiments
further cap each container's memory to force disk-backed execution
(\autoref{sec:eval:ooc} details the memory budgets).

%% --------------------------------------------------------------------------
\section{Datasets}
\label{sec:eval:datasets}

\paragraph{Structural graphs.}
\autoref{tab:datasets} lists the structural graphs used in the static and
dynamic evaluations.
They span three orders of magnitude in edge count (51\,M to 17\,B) and in
on-disk size (0.8\,GB to 339\,GB).
We generate the Graph500 and uniform random graphs with the Teseo graph
generator~\cite{teseo2021deLeo}; Dota~League, Twitter~MPI, and com-friendster
come from the LDBC Graphalytics benchmark
suite~\cite{LDBCGrapalytics22:online}; uk-2007 comes from the Koblenz Network
Collection~\cite{kunegis2013konect}.

%%%%%%%%%%%%%%%%%%%%%%%%%%%%%%%%%%%%%%%%%%%%%%%%%%%%%%%%%%%%%%%%%%%%%%%%%%%%%%%%
\begin{table}[h]
  % \footnotesize
  \begin{tabular*}{\columnwidth}{@{\extracolsep{\fill}}lrrrr@{}}
  \toprule
  Dataset Name &
    Vertices $|V|$ &
    Edges $|E|$ &
    Avg.\ Degree &
    Size (GB) \\ \midrule
  Dota-league             & 61.17\,K  & 50.87\,M  & 831.45 & 0.8    \\
  graph500-22, uniform-22 & 2.39\,M   & 64.15\,M  & 26.84  & 0.95   \\
  graph500-24, uniform-24 & 8.87\,M   & 260.38\,M & 29.36  & 4.34   \\
  graph500-26, uniform-26 & 32.80\,M  & 1.05\,B   & 32.01  & 18.59  \\
  Twitter MPI             & 52.58\,M  & 1.96\,B   & 37.28  & 29.63  \\
  com-friendster          & 65.60\,M  & 1.80\,B   & 27.44  & 32.36  \\
  uk-2007                 & 105.15\,M & 3.30\,B   & 31.38  & 54.10  \\
  Graph500-28             & 121.24\,M & 4.23\,B   & 34.89  & 75.64  \\
  Graph500-30             & 447.80\,M & 17.02\,B  & 38.00  & 339.25 \\ \bottomrule
  \end{tabular*}
  \caption{Structural graph datasets.
    Vertices $|V|$, edges $|E|$, average degree ($|E|/|V|$), and raw
    edge-list size on disk.}
  \label{tab:datasets}
  \end{table}
%%%%%%%%%%%%%%%%%%%%%%%%%%%%%%%%%%%%%%%%%%%%%%%%%%%%%%%%%%%%%%%%%%%%%%%%%%%%%%%%

\paragraph{Property graphs.}
The property graph evaluation uses five datasets beyond the structural graphs
listed above.
The LDBC Social Network Benchmark (SNB) at scale factor~3 defines a property
graph with ${\sim}800$\,K vertices and ${\sim}3.5$\,M edges across 4~entity
types (Person, Post, Comment, Forum) and 7~relation types (schema in
\autoref{fig:ldbc-snb-schema}).
We use cit-Patents (3.77\,M vertices, 16.5\,M edges) for graph algorithm
comparisons against graph databases.
The microbenchmarks use two structural graphs, DBLP co-authorship
(${\sim}317$\,K vertices, ${\sim}1$\,M edges) and WikiTalk (${\sim}2.4$\,M
vertices, ${\sim}5$\,M edges), and two property graphs from the AsterDB
artifact, LDBC and Freebase.

%% --------------------------------------------------------------------------
\section{Comparison Systems and Benchmark Suites}
\label{sec:eval:systems}

\autoref{tab:comparison-systems} lists all comparison systems organized by
evaluation context.
\autoref{sec:related-work} describes each system in detail; this section
identifies the benchmark framework for each context.

\paragraph{Static analytics.}
The in-memory (\autoref{sec:eval:inmem}) and out-of-core
(\autoref{sec:eval:ooc}) experiments use LDBC Graphalytics benchmark
kernels~\cite{LDBCGrapalytics22:online} (BFS, PageRank, connected components,
triangle counting, and SSSP), each implemented against the system's native
API\@.
GAPBS provides reference implementations on native CSR; Galois, Ligra, and
Gemini each provide their own implementations.
The out-of-core systems (GraphChi, X-Stream, Blaze) also implement these
kernels natively.

\paragraph{Dynamic graphs.}
The dynamic experiments (\autoref{sec:eval:dynamic}) use the GFE benchmarking
driver~\cite{teseo2021deLeo}, which exposes a uniform insertion and analytics
interface across Teseo, LiveGraph, GraphOne, and \project{}.

\paragraph{Property graphs.}
The property graph experiments (\autoref{sec:eval:property}) use the AsterDB
evaluation artifact~\cite{moAsterEnhancingLSMstructures2025}, a Docker-based
benchmark suite with Apache TinkerPop's Gremlin traversal language as a uniform
query interface.
Each graph database implements a TinkerPop \emph{provider} that translates
Gremlin traversals into its native storage operations.
\project{} bypasses this layer and runs benchmarks through its native C++ API;
we quantify the overhead of the Gremlin layer
in~\autoref{sec:eval:property-ldbc}.
The LDBC SNB queries are additionally compared against
NeuG~\cite{NeuG2025Alibaba}, an embedded mmap-based CSR system that accepts
Cypher queries.

%%%%%%%%%%%%%%%%%%%%%%%%%%%%%%%%%%%%%%%%%%%%%%%%%%%%%%%%%%%%%%%%%%%%%%%%%%%%%%%%
\begin{table}[t]
  \centering
  % \footnotesize
  \begin{tabular*}{\columnwidth}{@{\extracolsep{\fill}}lll@{}}
    \toprule
    Context & System & Key Characteristic \\
    \midrule
    \multirow{4}{*}{\shortstack[l]{In-memory\\analytics}}
    & GAPBS   & Native CSR reference implementations \\
    & Galois  & Priority-driven irregular parallelism \\
    & Ligra   & Direction-optimizing traversal \\
    & Gemini  & Chunk-based partitioning \\
    \midrule
    \multirow{3}{*}{\shortstack[l]{Out-of-core\\analytics}}
    & GraphChi  & Parallel sliding windows \\
    & X-Stream  & Edge-centric streaming \\
    & Blaze     & NVMe-targeted scatter-gather \\
    \midrule
    \multirow{3}{*}{\shortstack[l]{Dynamic\\graphs}}
    & Teseo     & Fat-tree, optimistic latching \\
    & LiveGraph & Transactional edge log, mmap-backed \\
    & GraphOne  & Lock-free ring buffer, background archiver \\
    \midrule
    \multirow{7}{*}{\shortstack[l]{Property\\graph}}
    & Neo4j      & Native graph database (Cypher) \\
    & ArangoDB   & Multi-model document store (AQL) \\
    & NeuG       & Embedded mmap-based CSR (Cypher) \\
    & AsterDB    & LSM-tree backed (native API) \\
    & JanusGraph & Distributed, pluggable backends \\
    & OrientDB   & Multi-model database \\
    & PostgreSQL & Relational via SQLG \\
    \bottomrule
  \end{tabular*}
  \caption{Comparison systems by evaluation context.}
  \label{tab:comparison-systems}
\end{table}
%%%%%%%%%%%%%%%%%%%%%%%%%%%%%%%%%%%%%%%%%%%%%%%%%%%%%%%%%%%%%%%%%%%%%%%%%%%%%%%%

%% --------------------------------------------------------------------------
\section{\project{} Configurations}
\label{sec:eval:configs}

\project{} supports two structural representations, Adjacency~List
(\autoref{arch:adjlist}) and EdgeKey (\autoref{arch:edgekey}), and two
property storage modes, embedded (\autoref{arch:prop-embedded}) and columnar
(\autoref{arch:prop-columnar}).
Each experiment uses a subset of the four combinations:

\begin{itemize}
  \item \textbf{Static analytics}: Adjacency~List only.
        EdgeKey stores each edge as an independent B-tree row; iterating every
        edge in the graph requires visiting each row individually, yielding
        scan performance orders of magnitude slower than the packed adjacency
        list format for full-graph iterative algorithms.

  \item \textbf{Dynamic graphs}: Both Adjacency~List and EdgeKey.
        EdgeKey's per-edge rows allow concurrent, contention-free insertion;
        Adjacency~List's packed rows deliver faster post-ingestion analytics.
        The tradeoff between ingest throughput and analytical performance is a
        central result of the dynamic evaluation
        (\autoref{sec:eval:dynamic}).

  \item \textbf{Property graph queries}: All four combinations
        (Adjacency~List and EdgeKey, each with embedded and columnar
        properties).
        The embedded-vs-columnar tradeoff is query-dependent: embedded
        mode favors multi-hop traversals that access properties at each hop;
        columnar mode favors analytical scans over single property columns.
        \autoref{sec:eval:emb-vs-col} quantifies this tradeoff.
\end{itemize}

%% --------------------------------------------------------------------------
\section{General Methodology}
\label{sec:eval:methodology}

Unless a specific evaluation section states otherwise:

\begin{itemize}
  \item All experiments run after warmup.
        We report the mean of five trials for static analytics, the mean of
        four warm trials for out-of-core experiments, and the median (p50)
        latency for property graph queries and microbenchmarks.
  \item We set timeouts at 2~hours for graph algorithm benchmarks.
  \item Static and dynamic experiments use all cores on one NUMA node
        (96~threads on Machine~A).
        Property graph queries run single-threaded to measure per-query latency.
  \item Each evaluation section describes its own experiment-specific
        methodology (memory budgets for out-of-core execution, query parameter
        curation for LDBC SNB, and batching strategies for client-server
        systems).
\end{itemize}
